%% Q1 ctx
\textbf{Directions for Q.~1 to 5:} Read the following passage and answer the questions.\par\smallskip
At first sight silence would seem to be the essence of non communication. Yet we know that in some circumstances silence can speak volumes. It is also a potent tool in the communication tool box, for the speaker, as it adds emphasis to point, allowing time for its full significance to sink in. Handled skillfully, it can build up anticipation for what is to come the pregnant pause. It gives opportunity to make eye contact with audience and gauge their reactions. It is an essential component of the timing of one's delivery. For the listener silence can buy time, either to digest what has been said or to consider a response, or both. Silence will, often lure the speaker into saying more than he or she may have intended.

%% Q1
The passage states directly: ``It is an essential component of \textbf{the timing of one's delivery}.''
\why{``Non-communication'' is what silence \emph{seems} to be at first sight, a view the passage then rejects. Gauging reactions is an \emph{opportunity} that silence gives. Distraction is not mentioned.}

%% Q2
``Handled skillfully, it can \textbf{build up anticipation} for what is to come: the pregnant pause.'' A pregnant pause is a deliberate silence full of meaning that makes the audience wait for what follows.

%% Q3
``For the listener silence can \textbf{buy time}, either to digest what has been said or to consider a response.''
\why{Emphasising a point, being a potent tool and speaking volumes are described as uses of silence for the \emph{speaker} or in general.}

%% Q4
``At first sight silence would seem to be the essence of \textbf{non communication}.''

%% Q5
``It [silence] gives opportunity to make \textbf{eye contact} with audience and gauge their reactions.''

%% Q6
Devices that both \emph{accept} and \emph{deliver} data are input/output devices:
\begin{itemize}
\item \textbf{Touch screen}: displays output and takes touch input.
\item \textbf{USB memory (pen drive)}: data can be written to it and read from it.
\item \textbf{CD drive} (read/write): reads from and writes to discs.
\end{itemize}
A \textbf{screen projector} only displays output. Hence a, b and c.

%% Q7
(a) True: cache is a small, very fast memory between the CPU and RAM that speeds up processing.
(b) True: primary memory (RAM/ROM) is also called main memory.
(c) True: the hard disk is secondary (auxiliary) storage.
(d) False: \textbf{SRAM is precisely the memory used for cache}, because it is fast and needs no refreshing. DRAM is used for main memory.
Hence a, b and c are true and d is false.

%% Q8
\textbf{CC BY-SA} (Creative Commons \emph{Attribution--ShareAlike}):
\begin{itemize}
\item allows copying, redistribution and publishing (a);
\item requires \textbf{attribution} to the creator (b);
\item permits commercial use, so (c) is wrong;
\item permits modification and adaptation, so (d) is wrong;
\item derivatives must be released \textbf{under the same licence}, so the licence cannot be changed (e).
\end{itemize}
The answer is a, b and e.
\tip{CC elements: BY = attribution, SA = share-alike, NC = non-commercial, ND = no derivatives. CC0 = public-domain dedication.}

%% Q9
The \emph{National Mission on Education through ICT (NMEICT)} lists as opportunities the availability of EduSAT, the abundance of knowledge on the internet, and the fast-growing network of cyber kiosks and cafes. The mission document looks for \textbf{low}-power, low-cost connectivity and devices (e.g.\ the low-cost \emph{Aakash} tablet), so ``high power consumption'' is \emph{not} an opportunity.

%% Q10
\textbf{SAKSHAT}, launched on 30 October 2006 by MHRD, was conceived as a \textbf{``One Stop Education Portal''} addressing the educational needs of learners from school to postgraduate level. It is the portal of the NMEICT.
\why{DIKSHA is the national platform for school teachers and students (2017); SAMARTH is an e-governance ERP for universities; Samagra Shiksha is a scheme for school education.}

%% Q11
The differences alternate between $-2$ and $-4$:
\[40\xrightarrow{-2}38\xrightarrow{-4}34\xrightarrow{-2}32\xrightarrow{-4}28\xrightarrow{-2}\mathbf{26}\]

%% Q12
Treat the three letter positions as three separate series (A=1, \ldots, Z=26, counting cyclically):
\begin{itemize}
\item 1st letters: A, D, H, M, S (+3, +4, +5, +6), so next $+7$ gives \textbf{Z}.
\item 2nd letters: B, G, M, T, B (+5, +6, +7, +8), so next $+9$: B(2)$+9=$ \textbf{K}.
\item 3rd letters: D, K, S, B, L (+7, +8, +9, +10), so next $+11$: L(12)$+11=$ \textbf{W}.
\end{itemize}
The next term is \textbf{ZKW}.

%% Q13
\begin{itemize}
\item First letters: C, F, I, \textbf{L} (+3).
\item Numbers: $4=2^2$, $9=3^2$, $16=4^2$, so next $5^2=$ \textbf{25}.
\item Last letters: X, U, R, \textbf{O} ($-3$).
\end{itemize}
The next term is \textbf{L25O}.

%% Q14
Decimal values: $\tfrac58=0.625$, $\tfrac{9}{13}\approx0.692$, $\tfrac{11}{17}\approx0.647$.
Descending order: $\tfrac{9}{13}>\tfrac{11}{17}>\tfrac58$, which is option (C).

%% Q15
\[\frac47\times\frac23\times\frac56\times\frac58\times1008=\frac{4\times2\times5\times5}{7\times3\times6\times8}\times1008=\frac{200}{1008}\times1008=\mathbf{200}\]

%% Q16
\textbf{Particulate matter}, especially PM$_{2.5}$, is regarded by WHO as the pollutant with the greatest health impact. The fine particles penetrate deep into the lungs and bloodstream and cause respiratory and cardiovascular disease and premature death.

%% Q17
UNESCO's \textbf{ESD for 2030} framework has five priority action areas: (1) advancing policy, (2) transforming learning environments, (3) building capacities of \emph{educators}, (4) empowering and mobilising youth, and (5) accelerating local-level actions. ``Building capacities of \textbf{administrators}'' is not one of them.

%% Q18 q
Match the following and choose the correct option given below.\\[2pt]
\begin{tabular}{@{}ll@{}}
\textbf{Movement} & \textbf{Activist}\\
I. Chipko Movement & a. M.~S.~Swaminathan\\
II. Narmada Bachao Movement & b. Medha Patkar\\
III. Green Revolution & c. Amrita Devi\\
IV. Bishnoi Movement & d. Sunderlal Bahuguna\\
\end{tabular}

%% Q18
\begin{itemize}
\item \textbf{Chipko} (1973, Uttarakhand): Sunderlal Bahuguna (also Chandi Prasad Bhatt and Gaura Devi), so I-d.
\item \textbf{Narmada Bachao Andolan} (1985): Medha Patkar, so II-b.
\item \textbf{Green Revolution} (1960s): M.~S.~Swaminathan, so III-a.
\item \textbf{Bishnoi movement} (1730, Khejarli, Rajasthan): Amrita Devi Bishnoi and 363 Bishnois gave their lives to save Khejri trees, so IV-c.
\end{itemize}

%% Q19
\emph{``The Future We Want''} is the outcome document of the \textbf{United Nations Conference on Sustainable Development (Rio+20)}, held in Rio de Janeiro in June 2012. It launched the process that led to the SDGs.
\why{The Brundtland Commission (World Commission on Environment and Development) produced \emph{Our Common Future} (1987).}

%% Q20
The \textbf{Nagoya Protocol} (adopted 2010, in force 2014), a supplementary agreement to the Convention on Biological Diversity, concerns \textbf{Access and Benefit-Sharing (ABS)}: the benefits arising from the use of genetic resources must be shared in a \textbf{fair and equitable} way.
\why{Option (D) describes the \emph{Cartagena Protocol on Biosafety} (2000), which covers living modified organisms.}

%% Q21 q
Study the chart and answer the question. The pie chart given here represents the domestic expenditure of a family in percent. If the total monthly income of the family is Rs.~55,660/-, after PF deductions and payment of the house rent, the total monthly income of the family remains

%% Q21 fig
\begin{tikzpicture}[scale=0.9]
\fill[brand!15] (0,0) -- (160:1.4) arc (160:235.6:1.4) -- cycle;
\draw (0,0) circle (1.4);
\draw (0,0) -- (160:1.4); \draw (0,0) -- (235.6:1.4);
\node[font=\scriptsize,align=center] at (198:0.85) {Child\\Educ.\\21\%};
\node[font=\scriptsize,align=left,right] at (1.7,0) {Other slices: PF, Food,\\House Rent, Entertainment,\\Medical (values not legible\\in the question booklet)};
\end{tikzpicture}

%% Q21
Only the \emph{Child Education = 21\%} slice was legibly printed. The PF and house-rent percentages, which are needed to answer, could not be read from the chart, so the question could not be solved and the agency \textbf{dropped} it.
For reference, the options are about 70\% of the income: $0.70\times55{,}660=38{,}962$ (option C). This would be the answer if PF and house rent together formed 30\% of expenditure.

%% Q22 opts
\tikz[baseline=3mm,scale=0.45]{\draw[->](0,0)--(2.6,0);\draw[->](0,0)--(0,1.6);\draw[thick,brand](0,1)--(1.8,1);}
\tikz[baseline=3mm,scale=0.45]{\draw[->](0,0)--(2.6,0);\draw[->](0,0)--(0,1.6);\draw[thick,brand](0,0.8) .. controls (1.2,0.8) and (1.8,1) .. (2,1.5);}
\tikz[baseline=3mm,scale=0.45]{\draw[->](0,0)--(2.6,0);\draw[->](0,0)--(0,1.6);\draw[thick,brand](0,0)--(1.9,1.5);}
\tikz[baseline=3mm,scale=0.45]{\draw[->](0,0)--(2.6,0);\draw[->](0,0)--(0,1.6);\draw[thick,brand](0,0.6) .. controls (0.4,1.3) and (0.6,1.4) .. (0.9,0.9) .. controls (1.1,0.6) and (1.3,0.9) .. (1.6,0.95) -- (2.3,0.95);}

%% Q22
The statement has three phases. Prices \textbf{rise} (January--February), then \textbf{fall slightly} (March), then stay \textbf{steady} (now). Only graph (D) shows a rise, a small dip and then a flat line.
\why{(A) is flat throughout; (B) and (C) keep rising, with (B) accelerating and (C) at a constant rate.}

%% Q23
Sampling interval $k=\dfrac{N}{n}=\dfrac{500}{20}=25$. Selected numbers: $9, 34, 59, 84, 109, 134, \ldots, 309, \ldots, 484$, i.e.\ $9+25m$.
\begin{itemize}
\item $134=9+25\times5$ \checkmark, $309=9+25\times12$ \checkmark, $484=9+25\times19$ \checkmark
\item $89-9=80$ is not a multiple of 25, so \textbf{89 is not selected}.
\end{itemize}

%% Q24 fig
\begin{tikzpicture}
\begin{axis}[ybar,width=10cm,height=4.4cm,bar width=11pt,ymin=0,ymax=5800,
  symbolic x coords={91-92,92-93,93-94,94-95,95-96,96-97,97-98,98-99},xtick=data,
  nodes near coords,every node near coord/.append style={font=\tiny},
  tick label style={font=\scriptsize},ylabel={\scriptsize million US\$},
  title={\scriptsize Foreign Exchange Reserve of a Country},axis x line*=bottom,axis y line*=left]
\addplot[fill=brand!70,draw=brand] coordinates {(91-92,2640) (92-93,3720) (93-94,2520) (94-95,3360) (95-96,3120) (96-97,4320) (97-98,5040) (98-99,3120)};
\draw[accent,dashed,thick] (axis cs:91-92,3480) -- (axis cs:98-99,3480);
\end{axis}
\end{tikzpicture}

%% Q24
\[\text{Average}=\frac{2640+3720+2520+3360+3120+4320+5040+3120}{8}=\frac{27840}{8}=3480\]
Above the average (dashed line): 3720, 4320 and 5040, which is \textbf{3 years}. Below: 2640, 2520, 3360, 3120 and 3120, which is \textbf{5 years}. The ratio is \textbf{3 : 5}.

%% Q25
\textbf{Data interpretation} means drawing meaning from analysed data, relating it to the research questions or hypotheses and \textbf{arriving at conclusions}. Classification, statistical analysis and comparison are \emph{earlier} steps that interpretation builds upon.

%% Q26
\textbf{Formative evaluation} (Michael Scriven, 1967) is carried out \emph{during} instruction to give continuous feedback, so that teaching can be adjusted and \textbf{students' learning improved}. It is assessment \emph{for} learning.
\why{Identifying learning difficulties is \emph{diagnostic} evaluation; promotion to the next class is based on \emph{summative} evaluation.}

%% Q27
\textbf{Shodhganga} (INFLIBNET) is a digital \emph{repository of Indian PhD theses}; it is not an online teaching platform.
\why{SWAYAM Prabha (educational DTH), MOOCs and Google Meet (live classes) are all used for online or distance teaching.}

%% Q28
True statements:
\begin{itemize}
\item every child is unique (1);
\item students of the same intelligence may achieve differently, because of motivation, study habits and so on (2);
\item heredity and environment together cause individual differences (3).
\end{itemize}
Statement 4 is wrong because differences are physical, social and emotional as well as psychological. Statement 5 is wrong because differences do affect learning and teaching. The answer is 1, 2 and 3.

%% Q29
A good teacher (i) finds out \emph{why} the student is quiet in class and talks to him or her, (ii) gives the student a topic to present so as to build confidence, and (iii) involves the student by asking questions. Leaving the student alone (4) or sending the student to play (5) neglects the teacher's responsibility. The answer is 1, 2 and 3.

%% Q30
Constructive handling of a bright student who disrupts class: set challenging work in the area of interest (2), suggest enrichment reading (3), and connect the student's points with the lesson so the contribution is valued (4).
Setting impossible work to humiliate the student (1) or complaining to the parents (5) is negative. The answer is 2, 3 and 4.

%% Q31
\textbf{Deep ecology} (Arne Naess, 1973) holds that all living beings have intrinsic value. It emphasises \textbf{ecosophy} (ecological wisdom or harmony), \textbf{richness and diversity of life}, and a \textbf{simple, sustainable lifestyle}. It is not limited to the abiotic ecosystem. The answer is a, b and c.

%% Q32
According to \textbf{AISHE 2019--20}, the \textbf{Gender Parity Index (GPI)} in higher education was \textbf{1.01}. This was the first time that female GER exceeded male GER. Total GER was 27.1\%.

%% Q33
\emph{Higher Education 4.0}, aligned with Industry 4.0, promotes technology-integrated, flexible, personalised learning. It stresses \textbf{blended learning}, combining online and face-to-face modes.

%% Q34 q
Match the following given in List--I and List--II and select the correct option given below.\\[2pt]
\begin{tabular}{@{}p{0.3\linewidth}p{0.62\linewidth}@{}}
\textbf{List--I} & \textbf{List--II}\\
1. National Digital Library & a. Plenty of international electronic journals and e-books are made available to all HEIs\\
2. National Academic Depository & b. To incorporate Robotics in Engineering Education\\
3. e-Yantra & c. To facilitate digital issuance, storage, access and verification of Academic credentials\\
4. e-Shodh Sindhu & d. To develop a framework of virtual repository of learning resources with a single window search facility\\
\end{tabular}

%% Q34 opts
1-d, 2-b, 3-a, 4-c
1-d, 2-a, 3-c, 4-b
1-b, 2-c, 3-d, 4-a
1-d, 2-c, 3-b, 4-a

%% Q34
\begin{itemize}
\item \textbf{National Digital Library of India} (IIT Kharagpur): a virtual repository of learning resources with single-window search, so 1-d.
\item \textbf{National Academic Depository (NAD)}: digital issue, storage, access and verification of academic awards, so 2-c.
\item \textbf{e-Yantra} (IIT Bombay): brings robotics into engineering education, so 3-b.
\item \textbf{e-ShodhSindhu} (INFLIBNET): e-journals and e-books for HEIs, so 4-a.
\end{itemize}

%% Q35
The fundamental principles in NEP~2020 include extensive use of technology, creativity and critical thinking, multidisciplinarity and a holistic approach (synergy in curriculum), and ``respect for diversity and respect for the local context''. \textbf{``Respect for global diversity''} is not among the stated principles.

%% Q36 q
Match Section--I with Section--II.\\[2pt]
\begin{tabular}{@{}ll@{}}
\textbf{Section--I} & \textbf{Section--II}\\
1. Haptics & a. Use of body movement\\
2. Proxemics & b. Use of space and effects\\
3. Paralanguage & c. Communication through touching\\
4. Kinesics & d. Use of pitch and volume of voice\\
\end{tabular}

%% Q36
Haptics is touch (1-c); Proxemics is space (2-b); Paralanguage covers pitch, volume, tone and rate of voice (3-d); Kinesics is body movement (4-a). The match is 1-c, 2-b, 3-d, 4-a.

%% Q37
Teleconferencing, voice mail and chat are \textbf{communication tools}: they carry messages between people, synchronously or asynchronously. A \textbf{multimedia encyclopaedia} is an information or reference resource, not a communication tool. The answer is 1, 2 and 3.

%% Q38
The \textbf{7 Cs of effective communication} are Clarity, Conciseness, Concreteness, Correctness, Coherence, Completeness and Courtesy (some versions add Consideration). \textbf{Cultural heterogeneity} is not one of them; cultural differences are in fact a potential \emph{barrier}.

%% Q39
The components of a communication system are \textbf{sender (source), message, encoding, channel, receiver, decoding and feedback}, along with noise and context. \textbf{Reinforcement} is a learning-theory concept (Skinner), not a component of communication.

%% Q40
\textbf{Semantic barriers} are barriers of \emph{meaning}. They include a badly expressed message (1), technical language or jargon (2), and unclarified or wrong assumptions (3). \textbf{Lack of attention} (4) is a \emph{psychological} barrier. The answer is 1, 2 and 3.

%% Q41
The Government has set up regulatory and quality-assurance bodies (UGC, AICTE, NAAC, NBA, NCTE and others) because some institutions fall short on quality. The Reason gives exactly why these bodies were established, so both are true and R explains A.

%% Q42
In Nyaya logic, \textbf{Vyapti} (invariable concomitance) is:
\begin{itemize}
\item the relation between the \emph{middle term} (hetu, e.g.\ smoke) and the \emph{major term} (sadhya, e.g.\ fire) (1);
\item literally a state of \emph{pervasion}, in which smoke is pervaded by fire (2);
\item the logical basis of inference and hence of the Nyaya syllogism, expressed in its \emph{Udaharana} step (3).
\end{itemize}
All three statements are true.

%% Q43
Letters are removed alternately from both ends and from the front:
\begin{center}
ATTRIBUTION $\xrightarrow{\text{drop 1st \& last}}$ TTRIBUTIO $\xrightarrow{\text{drop first 2}}$ RIBUTIO $\xrightarrow{\text{drop 1st \& last}}$ IBUTI $\xrightarrow{\text{drop first 2}}$ \textbf{UTI}
\end{center}

%% Q44
$M+N$: M is the father of N. $N+O$: N is the father of O. So M is the father of O's father, i.e.\ O's \textbf{grandfather}.

%% Q45 opts
\tikz[baseline=-0.5ex,scale=0.3]{\draw (0,0) circle(1.2); \draw (0,0) circle(0.6); \draw (2.2,0) ellipse(0.6 and 0.4);}
\tikz[baseline=-0.5ex,scale=0.3]{\draw (0,0.5) circle(0.7); \draw (1.3,0.5) circle(0.7); \draw (0.65,-0.6) circle(0.7);}
\tikz[baseline=-0.5ex,scale=0.3]{\draw (0,0) circle(1.3); \draw (-0.55,0) circle(0.5); \draw (0.55,0) circle(0.5);}
\tikz[baseline=-0.5ex,scale=0.3]{\draw (0,0) circle(0.8); \draw (1.1,0) circle(0.8); \draw (2.2,0) circle(0.8);}

%% Q45
In the official answer, \textbf{Travellers} is the large set, and the two means of travel, \textbf{train} (train travellers) and \textbf{bus} (bus travellers), are two separate groups inside it. They are shown as two disjoint circles, since a journey is made by one or the other, within one big circle. This is diagram (C).
\begin{center}
\begin{tikzpicture}[scale=0.7,font=\scriptsize]
\draw[fill=brand!8] (0,0) circle(1.6); \node at (0,1.25) {Travellers};
\draw[fill=white] (-0.7,-0.1) circle(0.6) node {Train}; \draw[fill=white] (0.7,-0.1) circle(0.6) node {Bus};
\end{tikzpicture}
\end{center}

%% Q46
An \textbf{achievement test} measures how well the content taught has been learnt, so it must have high \textbf{content validity} (it must sample the syllabus adequately). An \textbf{entrance test} is used to predict future success in a course, so it must have high \textbf{predictive validity}.

%% Q47
The sequence \emph{open coding $\rightarrow$ selective coding $\rightarrow$ theoretical coding} belongs to Glaser's version of \textbf{grounded theory}. Strauss and Corbin's version uses open $\rightarrow$ axial $\rightarrow$ selective coding. In both, theory is generated from data by constant comparison.

%% Q48
A reaction (rating) scale scored 0--5 has equal intervals between successive points. In educational research such summated rating data are treated as \textbf{interval-scale} data, which allows means and standard deviations to be computed. The official key accepts this answer.
\tip{Scales: Nominal (names or categories), Ordinal (rank order), Interval (equal intervals, arbitrary zero), Ratio (equal intervals with a true zero).}

%% Q49
A \textbf{normal distribution} is perfectly symmetrical (\textbf{skewness $=0$}, mean $=$ median $=$ mode) and has \textbf{mesokurtic} peakedness (kurtosis $=3$, or excess kurtosis $=0$).
\why{Leptokurtic curves are more peaked and platykurtic curves flatter than normal.}

%% Q50
Auguste Comte's \emph{Law of Three Stages} holds that human thought passes from the \textbf{theological} to the \textbf{metaphysical} and then to the \textbf{positive (scientific)} stage. This transformation to rational, scientific thinking gave rise to \textbf{Positivism}, the philosophy that only knowledge verified through observation and the scientific method is valid.
